\documentclass[11pt]{article}

\usepackage[margin=0.85in]{geometry}
\linespread{0.97}
\emergencystretch=3em
\usepackage{microtype}
\usepackage{amsmath}
\usepackage{booktabs}
\usepackage{array}
\usepackage{xcolor}
\usepackage{listings}
\usepackage[hidelinks]{hyperref}
\usepackage{titlesec}

\titleformat{\section}{\large\bfseries}{\thesection.}{0.6em}{}
\titleformat{\subsection}{\normalsize\bfseries}{\thesubsection}{0.6em}{}
\titlespacing*{\section}{0pt}{0.8ex plus 0.3ex minus 0.2ex}{0.3ex plus 0.1ex}
\titlespacing*{\subsection}{0pt}{0.5ex plus 0.2ex minus 0.1ex}{0.2ex plus 0.1ex}

\setlength{\parskip}{0.25ex plus 0.15ex}

\lstset{
  basicstyle=\ttfamily\small,
  keywordstyle=\color{blue!70!black}\bfseries,
  commentstyle=\color{gray}\itshape,
  stringstyle=\color{red!60!black},
  language=C++,
  showstringspaces=false,
  columns=fullflexible,
  breaklines=true,
  xleftmargin=1em,
  aboveskip=0.4ex,
  belowskip=0.4ex,
}

\title{\vspace{-2em}LLM-Driven Evolution of L2 Cache Prefetchers in ChampSim}
\author{15-740 Final Project Report}
\date{}

\begin{document}
\maketitle
\vspace{-3em}

\section*{0.\quad Abstract}

Hardware prefetcher design exposes a large hyper-parameter space:
classification thresholds, prefetch degrees, scoring formulas, and the
control-flow of the classifier itself. We apply OpenEvolve, an
LLM-in-the-loop evolutionary search system, to three production-grade L2
prefetchers in ChampSim~--- the Best-Offset Prefetcher (BOP), the
IP-Classifier Prefetcher (IPCP), and Berti~--- and use a $90{:}10$
IPC-versus-accuracy fitness function on \texttt{602.gcc\_s} as the
search signal. Twenty iterations per prefetcher discovers configurations
that lift gcc IPC by $27$--$34\%$ over the as-shipped versions and yield
combined scores of $2.08$ (BOP), $2.18$ (IPCP), and $2.20$ (Berti). The
IPCP gain comes from a structural mutation~--- making
\texttt{classify\_ip()} non-static so it can read
\texttt{global\_stream\_dir}~--- found at iteration~1 and never
exceeded. As an additional line of improvement we then port the
evaluator to a multi-trace fitness signal across an eight-trace SPEC
CPU~2017 subset and run a further 58 candidates on IPCP; the resulting
\texttt{IPCP-v40} configuration lifts the multi-trace combined score
from $1.243$ to $1.316$ ($+5.9\%$) and the geometric-mean speedup from
$1.293\times$ to $1.372\times$, with no per-trace regression. The
multi-trace finding is methodological rather than mechanical: the
single-trace optimum was a local optimum on gcc, and a single
hyper-parameter (\texttt{GS\_DEGREE}) carries the bulk of the headroom
the gcc-only signal could not see.

\section{Introduction}

Hardware prefetchers hide DRAM latency by predicting future memory
accesses. Modern designs expose dozens of numeric thresholds, scoring
functions, and classification rules; the design space is large and
sensitive to the workload mix used during tuning, and most published
configurations are hand-tuned by their authors against a small set of
benchmarks. This project asks whether LLM-driven evolutionary search can
discover comparable or better configurations automatically.

We use OpenEvolve, an open-source reimplementation of the AlphaEvolve
pattern: an LLM proposes diff-based mutations to a marked region of
source code, each candidate is compiled and scored by a user-supplied
evaluator, and a multi-island genetic algorithm maintains a population.
The marked region (\texttt{EVOLVE-BLOCK}) gives the LLM tightly scoped
write access to the parts of the prefetcher we are willing to vary, and
the rest of the simulator is untouched. We target three L2 prefetchers
(BOP, IPCP, Berti) and a single SPEC CPU~2017 trace, \texttt{602.gcc\_s},
with a $90{:}10$ IPC-vs-accuracy fitness function.

A natural question is whether single-trace fitness generalises. As an
additional line of improvement we replace the gcc-only evaluator with a
multi-trace fitness signal aggregated geometrically across an eight-trace
SPEC subset and re-run the IPCP evolution. The resulting v40
configuration both improves the multi-trace score by $5.9\%$ and
narrowly tops the single-trace gcc number.

\section{Related Work}

\textbf{IPCP}~(Pakalapati \& Panda, ISCA 2020) defines the four-class
IP-Classifier prefetcher (CS, CPLX, GS, NL) we evolve.
\textbf{Best-Offset Prefetcher}~(Michaud, HPCA 2016) selects a single
prefetch offset by scoring candidate offsets against a recent-access
filter. \textbf{Berti}~(Navarro Torres et al., MICRO 2022) is a
timing-aware delta prefetcher that ranks observed strides by how
\emph{timely} the implied prefetch would have been against the cache's
average miss latency. Reduced ports of all three are available in the
\texttt{ChampSim} repository (Gober et al., 2022) and we treat their
as-shipped configurations as the ``stock'' baselines.

\textbf{OpenEvolve} is an open-source reimplementation of the
\textbf{AlphaEvolve} pattern (Romera-Paredes et al., 2024). The LLM is
prompted with the current code, a system message describing the
fitness function, and the per-candidate score history; it returns a
diff that the framework applies, compiles, and re-scores. We use
\texttt{gpt-5-mini} via the CMU LiteLLM gateway with a population of
$30$~per island, two islands, $20$~iterations, an exploitation ratio of
$0.7$, and a maximum candidate length of $8000$ tokens.

\textbf{ChampSim} is the trace-driven simulator. Our configuration is a
$4$-GHz OoO core with a $352$-entry ROB, $64{\times}8$ L1I,
$64{\times}12$ L1D, $1024{\times}8$ L2C ($10$-cycle), $2048{\times}16$
LLC ($20$-cycle), and DDR4-3200 ($t_{\text{CAS}}{=}t_{\text{RCD}}{=}t_{\text{RP}}{=}24$).
The L2 prefetcher operates on physical block addresses; cross-page
prefetches are filtered out by \texttt{issue\_prefetch}'s page-boundary
check.

\section{Method}

\subsection{Evolution setup}

For each prefetcher we mark a contiguous region of the C++ source with
\texttt{// EVOLVE-BLOCK-START} / \texttt{// EVOLVE-BLOCK-END}. The
region contains the configurable constants, scoring functions, offset
lists, and classifier logic that we are willing to vary. The rest of the
prefetcher (state schema, hooks called by ChampSim) is fixed. The LLM
is given a system message describing the role of each constant and is
explicitly told what it may and may not change.

The fitness function on \texttt{602.gcc\_s} (200K warmup + 2M sim
instructions) is
\[
\textsf{score} \;=\; 0.9 \cdot \frac{\textsf{IPC}}{0.471} \,+\,
0.1 \cdot \frac{\textsf{useful}}{\textsf{useful}+\textsf{useless}}
\]
where $0.471$ is the no-prefetcher gcc IPC. A build or run failure
returns score $0$; the genetic algorithm's archive ensures the best
candidate is preserved across iterations. Each iteration is one LLM
call ($\approx 30$\,s) plus one ChampSim build ($\approx 60$\,s) plus
one simulation ($\approx 30$\,s).

\subsection{Multi-trace extension}

After the single-trace evolutions completed, we ported the evaluator to
an eight-trace SPEC CPU~2017 subset~--- perlbench\_s, gcc\_s, mcf\_s,
lbm\_s, omnetpp\_s, xalancbmk\_s, deepsjeng\_s, xz\_s (one region each,
${\approx}3.6$\,GiB total) spanning branch-heavy, memory-bound, and
streaming behaviour. The fitness function aggregates per-trace
speedups and accuracies as
\[
\mathrm{combined\_score}
\;=\; 0.9 \cdot \mathrm{geomean}_T(\mathrm{speedup}_T)
   \,+\, 0.1 \cdot \mathrm{mean}_T(\mathrm{accuracy}_T)
\]
where $\mathrm{speedup}_T = \mathrm{IPC}_T / \mathrm{IPC}^{\text{no-pf}}_T$
is computed against a per-trace no-prefetcher baseline cached in
\texttt{baselines\_ipcp\_multi.json}. The geometric mean prevents one
strong trace from masking weakness on another. The $0.9{:}0.1$ weighting
is unchanged so single-trace numbers under either evaluator are directly
comparable when restricted to gcc. We re-evaluated IPCP on the
multi-trace evaluator across $58$ candidates (\texttt{v1}--\texttt{v58})
to test whether additional headroom remained beyond the single-trace
optimum.

\section{Experimental Results}

\subsection{Single-trace results on \texttt{602.gcc\_s}}

\begin{center}\small
\begin{tabular}{l r r r r}
\toprule
Prefetcher & IPC & Speedup & Accuracy & Combined \\
\midrule
no prefetcher                & 0.471 & $1.00\times$ & ---     & ---  \\
BOP   (stock)                & 0.821 & $1.74\times$ & $98.7\%$ & 1.66 \\
IPCP  (stock)                & 0.811 & $1.72\times$ & $98.7\%$ & 1.64 \\
Berti (stock)                & 0.956 & $2.03\times$ & $96.8\%$ & 1.92 \\
\midrule
\textbf{BOP   (evolved, iter 6)}   & \textbf{1.040} & $\mathbf{2.21\times}$ & $96.5\%$ & \textbf{2.08} \\
\textbf{IPCP  (evolved, iter 1)}   & \textbf{1.088} & $\mathbf{2.31\times}$ & $97.4\%$ & \textbf{2.18} \\
\textbf{Berti (evolved, iter 15)}  & \textbf{1.100} & $\mathbf{2.34\times}$ & $94.7\%$ & \textbf{2.20} \\
\bottomrule
\end{tabular}
\end{center}

All three prefetchers improve on their hand-tuned stock configurations.
The relative IPC gains over stock are BOP $+27\%$, IPCP $+34\%$, and
Berti $+15\%$. Single-trace evolution saturates quickly: BOP plateaued at
iteration~$6$ of $20$, IPCP at iteration~$1$, and Berti at
iteration~$15$.

\textbf{What the LLM changed.} For BOP, the mutation doubled
\texttt{PREFETCH\_DEGREE} ($1{\to}2$), tightened
\texttt{MSHR\_THRESHOLD} ($0.7{\to}0.6$), shortened
\texttt{ROUND\_MAX} ($100{\to}64$), and duplicated the small offsets
inside \texttt{OFFSET\_LIST} so good candidates received multiple
independent trials. For Berti, the mutation rebalanced the scoring
weights (\texttt{TIMELINESS\_WEIGHT}~$1.0{\to}1.8$,
\texttt{ACCURACY\_WEIGHT}~$0.5{\to}1.6$), introduced a non-linear
late-penalty $\textsf{late\_ratio}^{1.5}$, and added a consistency
bonus that scales with $\sqrt{\text{coverage}}$.

The IPCP gain is the most interesting because it is structural rather
than numeric. Stock IPCP declares \texttt{classify\_ip()} as a static
member function, so the classifier sees only its four explicit
arguments (the IP-local stride history) and cannot read the workload-wide
\texttt{global\_stream\_dir} signal that
\texttt{prefetcher\_cache\_operate} maintains one scope up. The
iteration-1 mutation removed the \texttt{static} keyword on both the
declaration and the definition and added a stream-intercept block:

\begin{lstlisting}
if (global_stream_dir != 0) {
  if ((global_stream_dir > 0 && new_stride > 0) ||
      (global_stream_dir < 0 && new_stride < 0)) {
    if (confidence < CS_CONFIDENCE_THRESHOLD || old_class == GS)
      return GS;
  }
}
\end{lstlisting}

This routes IPs that happen to be moving in the global workload's
direction~--- but have not yet accumulated the confidence threshold for
the CS class~--- straight to the GS class, which issues a 3-deep
$\pm 1$ prefetch chain. On \texttt{gcc\_s} this re-classifies a large
population of ``going-with-the-flow'' PCs from NL (one prefetch) to GS
(three prefetches) and lifts IPC from $0.811$ to $1.088$.

\subsection{Multi-trace extension: IPCP-v40}

The single-trace results above leave open whether the IPC gains
generalise. We re-ran the evolved BOP, IPCP, and Berti binaries on the
eight-trace SPEC subset against a no-prefetcher baseline:

\begin{center}\small
\begin{tabular}{l r r r}
\toprule
Prefetcher & Geomean speedup & Mean accuracy & Combined \\
\midrule
no prefetcher                       & 1.000  & ---     & ---   \\
IPCP (stock)                        & 1.195  & $82.5\%$ & 1.158 \\
BOP   (evolved)                     & 1.218  & $82.2\%$ & 1.180 \\
IPCP  (evolved, iter 1)             & 1.293  & $79.1\%$ & 1.243 \\
Berti (evolved, iter 15)            & 1.294  & $78.0\%$ & 1.243 \\
\textbf{IPCP-v40 (multi-trace)}     & \textbf{1.372}  & \textbf{$81.0\%$} & \textbf{1.316} \\
\bottomrule
\end{tabular}
\end{center}

The single-trace-evolved BOP, IPCP, and Berti retain meaningful gains
across the multi-trace suite~--- IPCP's geomean speedup is $1.293\times$
and Berti's is $1.294\times$~--- but their combined scores cluster
around $1.24$, far below their gcc-only $2.0+$ figures. Both are
catastrophically inaccurate on memory-bound traces: IPCP issues 168K
prefetches on \texttt{mcf} with only $28.6\%$ accuracy, and Berti issues
166K with $20.9\%$ accuracy. As an additional line of improvement we ran
a further $58$ candidates on the multi-trace evaluator. The champion
(\texttt{IPCP-v40}) raises the IPCP combined score from $1.243$ to
$1.316$ ($+5.9\%$) and the geomean speedup to $1.372\times$.

The per-trace breakdown shows where the gain comes from:

\begin{center}\small
\begin{tabular}{l r r r r r r}
\toprule
Trace & no-pf & stock IPCP & IPCP iter-1 & \textbf{v40} & speedup vs.\ no-pf & $\Delta$ IPC vs iter-1 \\
\midrule
\texttt{600.perlbench\_s} & 1.935 & 1.986 & 1.995 & 1.997 & $1.03\times$ & $+0.002$ \\
\texttt{602.gcc\_s}       & 0.469 & 0.806 & 1.088 & 1.106 & $2.36\times$ & $+0.018$ \\
\texttt{605.mcf\_s}       & 0.210 & 0.270 & 0.260 & 0.274 & $1.31\times$ & $+0.014$ \\
\texttt{619.lbm\_s}       & 0.583 & 0.615 & 0.708 & 0.758 & $1.30\times$ & $+0.050$ \\
\texttt{620.omnetpp\_s}   & 0.226 & 0.232 & 0.230 & 0.232 & $1.03\times$ & $+0.002$ \\
\textbf{\texttt{623.xalancbmk\_s}} & 0.411 & 0.585 & 0.739 & \textbf{1.024} & $\mathbf{2.49\times}$ & $\mathbf{+0.285}$ \\
\texttt{631.deepsjeng\_s} & 1.090 & 1.106 & 1.109 & 1.112 & $1.02\times$ & $+0.003$ \\
\texttt{657.xz\_s}        & 2.156 & 2.508 & 2.508 & 2.508 & $1.16\times$ & $+0.000$ \\
\bottomrule
\end{tabular}
\end{center}

\texttt{xalancbmk} dominates the gain ($+0.285$ IPC, $+38\%$ over
iter-1 IPCP), followed by \texttt{lbm} ($+0.05$). On \texttt{gcc} v40
also slightly exceeds the iter-1 IPCP IPC ($1.106$ vs.~$1.088$), which
means the multi-trace optimum is also a slightly better gcc-only
optimum (combined $2.21$ vs.~$2.18$). On \texttt{mcf} the iter-1 IPCP
actually regresses below stock ($0.260$ vs.~$0.270$): the stream-aware
classifier introduced 60K extra useless prefetches there. v40's new
\texttt{GS\_STRIDE\_LIMIT}$=8$ gate~--- a one-line addition to
\texttt{classify\_ip()}~--- recovers the lost \texttt{mcf} accuracy
while keeping all of the streaming-trace gains.

\textbf{What the multi-trace evolution changed.} v40 differs from the
iter-1 IPCP in five constants and one classifier line:
\begin{lstlisting}
constexpr int CS_DEGREE       = 8;     // was 4
constexpr int CPLX_DEGREE     = 8;     // was 3
constexpr int GS_DEGREE       = 16;    // was 3 -- the dominant change
constexpr int GS_STRIDE_LIMIT = 8;     // new constant
// CS_CONFIDENCE_THRESHOLD, CPLX_CONFIDENCE_THRESHOLD,
// STREAM_DETECT_THRESHOLD, CONFIDENCE_SAT_MAX, MSHR_THRESHOLD: unchanged

if (dir_match && std::abs(new_stride) <= GS_STRIDE_LIMIT) {
  if (confidence < CS_CONFIDENCE_THRESHOLD || old_class == GS) return GS;
}
\end{lstlisting}
\texttt{GS\_DEGREE} carries the bulk of the gain ($3{\to}16$); the
limit is reached because \texttt{xalancbmk} absorbs prefetch chains all
the way to the 4-KiB page boundary, where \texttt{issue\_prefetch}'s
same-page check silently caps any further degree increase.

\subsection{Negative results around v40}

We invested $18$ further iterations probing structural changes after
v40. None beat it by more than measurement noise, and three exposed
instructive failure modes. Routing GS prefetches at the IP's actual
stride instead of $\pm 1$ (\texttt{v50}) lost $2.6\%$~--- the $\pm 1$
chain captures the immediately-adjacent cache line which is hot under
spatial locality. Removing the page-boundary check inside
\texttt{issue\_prefetch} (\texttt{v54}) dropped \texttt{lbm} accuracy
from $100\%$ to $72\%$ and regressed combined score to $1.281$;
ChampSim does propagate cross-page prefetches but every one is marked
useless. Boosting confidence by $+2$ on \texttt{useful\_prefetch=true}
(\texttt{v55}) overshot CS classifications and regressed to $1.278$.
These confirm v40 saturates the local optimum reachable with IPCP's
existing state schema and \texttt{EVOLVE-BLOCK} surface.

\section{Goals and Next Steps}

The proposal goal was to demonstrate that LLM-driven evolutionary
search can match or beat hand-tuned prefetcher configurations, and to
characterise the regime in which it works. We met that goal on the
single-trace metric for all three prefetchers ($+27$ to $+34\%$ IPC
over hand-tuned), and the multi-trace extension shows that an
additional $+5.9\%$ combined score is recoverable on IPCP once the
fitness signal is the right one.

Three directions are the most promising next steps. First, the
page-boundary plateau invites widening the evolvable scope to include
\texttt{prefetcher\_cache\_operate} and \texttt{issue\_prefetch}; with
that surface in scope, an evolved prefetcher could compute
remaining-lines-in-page dynamically and reallocate the saved degree
elsewhere. Second, the \texttt{useful\_prefetch} and \texttt{cache\_hit}
signals are entirely unused by the current IPCP; a more careful
integration~--- tracking a per-IP rolling accuracy in the IP table
rather than perturbing global confidence~--- could unlock the
closed-loop feedback we attempted in \texttt{v55}. Third, applying the
multi-trace evaluator to BOP and Berti would test whether the
single-hyper-parameter pattern (a single dominant degree) generalises
beyond IPCP; both currently sit below v40 on the multi-trace suite.

\section{Collaboration}

\section{Conclusion}

LLM-driven evolutionary search discovers prefetcher configurations
competitive with hand tuning. Twenty iterations against a single SPEC
trace lifted gcc IPC by $27$--$34\%$ on three prefetchers and surfaced
one structurally meaningful mutation (the non-static
\texttt{classify\_ip()} on IPCP). Re-evaluating those configurations on
a multi-trace SPEC subset shows that the single-trace gains do generalise
in geomean speedup ($1.29$--$1.31\times$) but that an additional
$+5.9\%$ combined score is available on IPCP once the fitness signal is
multi-trace. The multi-trace champion (\texttt{IPCP-v40}) raises geomean
speedup to $1.372\times$ via a single dominant hyper-parameter
(\texttt{GS\_DEGREE}~$3{\to}16$) and a small magnitude-aware gate that
prevents the more aggressive prefetcher from polluting memory-bound
traces. The plateau at v40 is enforced by ChampSim's page-boundary
check and is therefore an architectural rather than algorithmic limit;
further progress requires widening the evolvable region or rethinking
IPCP's state schema.

\end{document}
